% Fragmento público generado desde propietarios íntegros.
% Residencia: 52.13.3.
% Fuente: ../incoming/radion_monografia_20260827/manuscrito/sections/13_sintesis.tex:185-247
\subsubsection{Contrôles négatifs du système radional}

La construction est réfutable en des points concrets :
\begin{enumerate}
\item Si \(S_E\notin(1,2)\), le quotient APP n'est pas un et la clôture change de
cellule.
\item Si l'incidence active ne produit pas \(N_6=90\), ou s'il n'y a pas deux feuillets,
\(20/81\) ne s'ensuit plus.
\item Si le vérificateur reçoit \(180/\pi\) ou \(\epsR\) avant de construire les
opérandes, la génération est circulaire.
\item Si la soustraction par triades ne stabilise pas des préfixes compatibles, la
section coinductive proposée n'existe pas.
\item Si \(|C^*|\ge A\), l'ellipse positive cesse d'exister.
\item Si \(a_Sb_S\ne2\hret\), le théorème d'aire échoue.
\item Si \(C^*=0\), on doit avoir \(d=0\), \(\eta=0\), \(\tau=i\) et
\(j=1728\).
\item Si l'aire de Klein du sous-disque ne suit pas \eqref{eq:klein-area-R}, la
thèse de profondeur infinie échoue.
\item Si \(F_{\mathrm{spec}}=\epsR^\circ\) est affirmé sans contre-terme indépendant, la
différence numérique \eqref{eq:chiR} le réfute.
\item La production de \(12:-1\) exige de conserver, avec les cardinaux
\(90/120\), les douze secteurs dodécaphasiques, l'unique couture orientée et
l'action du lecteur sur ces générateurs. Changer une multiplicité
doit modifier l'image \(12S_{90}-S_{120}\). Pour la chaîne fondamentale
sans altération, le coefficient de \((1-q)^{-1}\) doit donner \(1/24\).
\item Si la fonction \(j\) est utilisée pour sélectionner le signe de charge, on perd
l'invariance \eqref{eq:j-invariance} et l'on commet une erreur de type.
\end{enumerate}

\subsubsection{Réalisation physique conjointe du radion}

L'image finale n'est pas un cercle corrigé. C'est une architecture de
publications :
\[
\begin{aligned}
&\text{APP:}&&\text{unité, reste, retenue et fibre};\\
&\text{TRIT:}&&\text{régime, orientation et feuillets};\\
&\text{TPK:}&&\text{phase, retour, dodécaphase et mémoire};\\
&\text{ellipse :}&&\text{action et anisotropie};\\
&\text{Klein:}&&\text{profondeur intérieure};\\
&\text{hélice :}&&\text{biographie du retour};\\
&\text{MD:}&&\text{horloge, énergie, charge, champ et causalité}.
\end{aligned}
\]

L'équation du radion réunit la face visible :
\BigRadionEquation
L'ellipse réunit la face d'action :
\[
\mathcal A(1)=\hret,\qquad \mathcal A(2\pi)=\hfull.
\]
La métrique de Klein réunit la profondeur :
\[
K=-1,\qquad \operatorname{Area}_K=\infty.
\]
L'holonomie réunit la mémoire :
\[
g(k+9)=g(k),\qquad B_{k+9}\ne B_k.
\]

\begin{thesisbox}
Le cercle est l'ombre angulaire ; l'ellipse est le corps d'action ; Klein mesure
l'intérieur ; l'hélice conserve la biographie. Le radion est l'unité qui
permet de lire les quatre affirmations sur un unique état sans réduire l'une à
l'autre.
\end{thesisbox}
% Fuente: ../incoming/radion_monografia_20260827/manuscrito/sections/B_valores_y_verificacion.tex:1-105
\subsubsection{Valeurs canoniques, stabilité et vérification reproductible}

\paragraph{Valeurs des invariants internes générés}

\begin{longtable}{p{.27\textwidth} p{.61\textwidth}}
\toprule
Objet & Valeur utilisée dans le profil coinductif\\
\midrule
\endfirsthead
\toprule
Objet & Valeur utilisée dans le profil coinductif\\
\midrule
\endhead
\(\pih\) & \(\begin{aligned}3.1415926535897932384626433832795028841\\[-2pt]971693993751058209749445923\ldots\end{aligned}\)\\
\(e_{\HMT}\) & évaluation coinductive du lecteur d'Euler, sans table externe\\
\(\alphah\) & \(\begin{aligned}0.0072973525692838009972851054723806629\\[-2pt]995641725396427091854046825\ldots\end{aligned}\)\\
\(A\) & \(\begin{aligned}7.2973525692838009972851054723806629995\\[-2pt]641725396427\ldots{}^\circ\end{aligned}\)\\
\(C^*\) & \(2.1237383389686457242619513803933607937\ldots{}^\circ\)\\
\(\eta\) & \(0.299689683921630799684926562863927\ldots\)\\
\(\hret\) & \(1.05457181691503906113369058472430\ldots\times10^{-34}U_S\)\\
\bottomrule
\end{longtable}

\paragraph{Convergence quantitative selon la profondeur}

\begin{center}
\small
\begin{tabular}{r l l}
\toprule
Profondeur & \(\epsR^\circ\) & lecture\\
\midrule
12 & \(5.1681282186551375597\times10^{-8}\) & préasymptotique\\
18 & \(5.1383017764316644535\times10^{-8}\) & 7 cifras estables\\
24 & \(5.1383016758575394560\times10^{-8}\) & 14 cifras estables\\
30 & \(5.1383016758575282072\times10^{-8}\) & 20 cifras estables\\
36 & \(5.138301675857528207168517\times10^{-8}\) & perfil largo\\
45 & \(5.13830167585752820716851646226459\times10^{-8}\) & largeur \(<4.81\times10^{-130}\)\\
\bottomrule
\end{tabular}
\end{center}

Les retours de profondeur \(9,18,27,36,45\) conservent la phase \(8\) et
augmentent la mémoire \(0,1,2,3,4\). Le rapport entre largeurs consécutives tend
vers
\[
3^{-54}=1.7196982334549856127610399316004871\ldots\times10^{-26}.
\]

\paragraph{Analyse de sensibilité structurelle par suppression contrôlée}

\begin{longtable}{p{.46\textwidth} p{.34\textwidth}}
\toprule
Opération & Sortie ou différence unitaire\\
\midrule
\endfirsthead
\toprule
Opération & Sortie ou différence unitaire\\
\midrule
\endhead
Opérateur canonique & \(8.96802822044562981999\times10^{-10}\)\\
Un feuillet, \(\rho=90/729\) & \(-1.3727056615960678\times10^{-5}\)\\
Phase \(\theta_1=20^\circ\) & \(+5.2359877649510169\times10^{-1}\)\\
Phase \(\theta_3=80^\circ\) & \(-5.2359877470149605\times10^{-1}\)\\
Sans canal \(e\) & \(1.8065350748904653\times10^{-3}\)\\
Torsion historique \(\Delta_4(A,C^*)\) & \(1.0517084608768816\times10^{-9}\)\\
Opérateur harmonique \(M_4\) moins canonique & \(1.0525500222895070\times10^{-17}\)\\
Queue sexadique moins canonique & \(2.0283936357433839\times10^{-12}\)\\
\bottomrule
\end{longtable}

La séparation des échelles montre une sensibilité structurelle : la clôture dépend
de la phase, des deux feuillets, des canaux \(\pi/e\), de la torsion globale et de la retenue.
Aucune ablation ne reproduit la valeur.

\paragraph{Certificats informatiques reproductibles}

L'ouvrage utilise les propriétaires exécutables suivants :
\begin{enumerate}
\item \sourcepath{output/RADION_TWOWAY_NONADICO_20260827/verificar_emergencia_radion_por_profundidad.py};
\item \sourcepath{output/RADION_TWOWAY_NONADICO_20260827/verificar_radion_helice_critica_tpk.py};
\item \sourcepath{output/RADION_TWOWAY_NONADICO_20260827/verificar_operador_tpk_target_free_profundidad.py};
\item le vérificateur consolidé de cette monographie, situé dans
\sourcepath{certificados/verificar_monografia_radion.py}.
\end{enumerate}
Les sorties attendues comprennent
\[
\texttt{PASS\_EMERGENCIA\_RADION\_POR\_PROFUNDIDAD},
\]
\[
\texttt{PASS\_RADION\_HELICE\_CRITICA\_TPK},
\qquad
\texttt{PASS\_MONOGRAFIA\_RADION\_HMT\_MD}.
\]

\paragraph{Distinction entre précision informatique et exactitude mathématique}

Qu'un programme affiche zéro à 120 ou 200 chiffres ne transforme pas à lui seul une
identité résiduelle en prédiction. L'exactitude mathématique procède de
\eqref{eq:cierre-unitario} ; la précision arbitraire permet d'évaluer ses
opérandes et de vérifier la convergence. La non-circularité procède de l'ordre
de construction et de l'interdiction des cibles, non du nombre de chiffres.

Le calcul Decimal peut laisser des résidus d'ordre \(10^{-218}\) lorsqu'il combine
des branches évaluées avec des précisions différentes. Ce nombre est un arrondi de
l'implémentation ; non une correction physique supplémentaire.
% Fuente: ../incoming/radion_monografia_20260827/manuscrito/sections/D_historia_fuentes.tex:1-169
\subsubsection{Généalogie conceptuelle et sources primaires}

\paragraph{Quantum d'action de Planck}

Planck introduit \(h\) dans la loi de rayonnement comme constante reliant énergie
et fréquence. La donnée historique décisive pour cet ouvrage est dimensionnelle :
\(h\) est une action. Le radion ne modifie pas ce rôle ; il explique pourquoi l'action
d'un tour et l'action par unité de phase sont liées par \(2\pi\).

\paragraph{Quantification angulaire de Bohr et de Sommerfeld}

Bohr utilise \(h/2\pi\) dans la quantification du moment cinétique. Sommerfeld étend
la quantification aux intégrales d'action et aux orbites képlériennes. Il n'est donc pas
historiquement exact de dire que Dirac a inventé la constante réduite. Sa
contribution appartient à une généalogie dans laquelle l'action angulaire était déjà
présente.

La nouveauté HMT--MD est différente : elle ne prend pas \(h/2\pi\) comme un quotient externe,
mais construit une ellipse dont le tour a pour aire \(h\) et dont le secteur d'un
radian a pour aire \(\hbar\).

\paragraph{Formalisme matriciel de Born et de Jordan et oscillateur quantique}

Born et Jordan placent le principe variationnel, les phases et l'oscillateur
harmonique au cœur de la mécanique matricielle ; le champ électromagnétique se
décompose en une collection d'oscillateurs. La réalisation LC du radion
prolonge cette intuition avec une donnée supplémentaire : le quotient de forme
\(e^{-\eta}\) détermine l'impédance relative tandis que produit, fréquence
et action restent fixes.

\paragraph{Fréquence angulaire et action dans la formulation de Dirac}

En 1925, Dirac écrit les fréquences en radians par unité de temps et utilise
des relations de la forme
\[
xy-yx=\frac{ih}{2\pi}[x,y],
\qquad
\Delta E=\frac{h}{2\pi}\omega.
\]
L'équation \(E=\hret\omega\) relie explicitement action réduite et
fréquence angulaire. Dans le radion, cette relation reçoit une géométrie : \(\hret\)
est l'aire balayée par une unité de phase.

\paragraph{Relations d'indétermination de Heisenberg, Robertson et Schrödinger}

Heisenberg formule la limitation conjointe des variables canoniquement
conjuguées. Robertson généralise l'inégalité ; Schrödinger intègre la
covariance. La matrice \eqref{eq:covariance} relève plus directement de cette
prolongation Robertson--Schrödinger que de l'article de Heisenberg pris isolément.

Aucun de ces auteurs ne formule l'ellipse généalogique complète de cet ouvrage.
L'affirmation historiquement défendable est qu'ils fournissent le langage dans
lequel l'ellipse d'action peut être reconnue comme un contour minimal ; la
construction du couple \((A,C^*)\), les foyers, Klein, l'hélice et la retenue constituent la
contribution HMT--MD.

\paragraph{Fonctionnelle de phase de Feynman}

La formulation spatio-temporelle attribue à chaque chemin une phase
\(e^{iS/\hbar}\). L'unité de radion rend la correspondance littérale :
incrémenter l'action de \(\hbar\) incrémente la phase d'un radian. La somme des
chemins reconnaît, dans un langage intégral, le même principe que l'ellipse
exprime comme aire.

\paragraph{Propagations avancée et retardée chez Wheeler--Feynman}

La théorie de l'absorbeur construit une interaction classique temporellement
symétrique au moyen de combinaisons de champs retardés et avancés. L'identité
HMT \(CU(t)C^{-1}=U(-t)\) et l'inversion des chemins fournissent un
antécédent opératoire. La promotion physique exige que l'application entre chemins
et propagateurs préserve action, orientation, frontière et forme symplectique.

La théorie de l'absorbeur n'est pas la même affirmation que l'interprétation de
Feynman et de Stückelberg des antiparticules comme propagation temporellement
inversée. Les deux peuvent reconnaître la bidirectionnalité ; leurs objets physiques et
leurs conditions de frontière sont différents.

\paragraph{Dualité de Hodge et séparation d'avec la cohomologie de Hochschild}

L'opération pertinente dans ce manuscrit est l'étoile de Hodge
\(\star^2=-I\) sur les deux-formes en signature lorentzienne. On n'a pas localisé dans
le corpus actif du radion de propriétaire d'homologie de Hochschild qui
agisse sur l'ellipse, la retenue ou le secteur \(90/120\). Introduire Hochschild
par ressemblance nominale dégraderait le typage. Si une prolongation future
construit un cocycle de déformation avec ses propres domaine et codomaine, il devra
entrer comme résultat nouveau, non comme rétroattribution.

\paragraph{Bibliographie primaire sélectionnée}

\begin{enumerate}
\renewcommand{\labelenumi}{[\arabic{enumi}]}
\item\hypertarget{radionbib:Planck1901}{}
M.~Planck,
\emph{Ueber das Gesetz der Energieverteilung im Normalspectrum},
Annalen der Physik 309 (1901), 553--563.
\url{https://doi.org/10.1002/andp.19013090310}

\item\hypertarget{radionbib:Bohr1913}{}
N.~Bohr,
\emph{On the Constitution of Atoms and Molecules},
Philosophical Magazine 26 (1913), 1--25.
\url{https://doi.org/10.1080/14786441308634955}

\item\hypertarget{radionbib:Sommerfeld1916}{}
A.~Sommerfeld,
\emph{Zur Quantentheorie der Spektrallinien, I},
Annalen der Physik 356 (1916), 1--94.
\url{https://doi.org/10.1002/andp.19163561702}

\item\hypertarget{radionbib:BornJordan1925}{}
M.~Born et P.~Jordan,
\emph{Zur Quantenmechanik},
Zeitschrift für Physik 34 (1925), 858--888.
\url{https://doi.org/10.1007/BF01328531}

\item\hypertarget{radionbib:Dirac1925}{}
P.~A.~M.~Dirac,
\emph{The Fundamental Equations of Quantum Mechanics},
Proceedings of the Royal Society A 109 (1925), 642--653.
\url{https://doi.org/10.1098/rspa.1925.0150}

\item\hypertarget{radionbib:Heisenberg1927}{}
W.~Heisenberg,
\emph{Über den anschaulichen Inhalt der quantentheoretischen Kinematik und Mechanik},
Zeitschrift für Physik 43 (1927), 172--198.
\url{https://doi.org/10.1007/BF01397280}

\item\hypertarget{radionbib:Dirac1928}{}
P.~A.~M.~Dirac,
\emph{The Quantum Theory of the Electron},
Proceedings of the Royal Society A 117 (1928), 610--624.
\url{https://doi.org/10.1098/rspa.1928.0023}

\item\hypertarget{radionbib:Robertson1929}{}
H.~P.~Robertson,
\emph{The Uncertainty Principle},
Physical Review 34 (1929), 163--164.
\url{https://doi.org/10.1103/PhysRev.34.163}

\item\hypertarget{radionbib:Schrodinger1930}{}
E.~Schrödinger,
\emph{Zum Heisenbergschen Unschärfeprinzip},
Sitzungsberichte der Preussischen Akademie der Wissenschaften (1930), 296--303.

\item\hypertarget{radionbib:WheelerFeynman1945}{}
J.~A.~Wheeler et R.~P.~Feynman,
\emph{Interaction with the Absorber as the Mechanism of Radiation},
Reviews of Modern Physics 17 (1945), 157--181.
\url{https://doi.org/10.1103/RevModPhys.17.157}

\item\hypertarget{radionbib:Feynman1948}{}
R.~P.~Feynman,
\emph{Space-Time Approach to Non-Relativistic Quantum Mechanics},
Reviews of Modern Physics 20 (1948), 367--387.
\url{https://doi.org/10.1103/RevModPhys.20.367}

\item\hypertarget{radionbib:Feynman1949}{}
R.~P.~Feynman,
\emph{The Theory of Positrons},
Physical Review 76 (1949), 749--759.
\url{https://doi.org/10.1103/PhysRev.76.749}
\end{enumerate}

\begin{narrativebox}
Les fondateurs contiennent les éléments : action, fréquence angulaire, variables
conjuguées, oscillateur, incertitude, phase des chemins et réversion temporelle.
La contribution du radion consiste à les réunir dans un objet généré : une ellipse
d'aire \(h\), un secteur d'aire \(\hbar\), un intérieur de profondeur infinie,
une orientation de feuillet et une retenue exacte de compatibilité.
\end{narrativebox}
% Fuente: ../incoming/radion_monografia_20260827/manuscrito/sections/E_conclusion.tex:1-107
\subsubsection{Synthèse finale : unité angulaire avec généalogie de phase, d'action et de mémoire}

Le radian ne disparaît pas. Il devient entièrement visible.

Dans son usage ordinaire, un radian est un rapport qui n'a pas besoin de se souvenir de la façon dont il a été
produit. L'arc et le rayon s'annulent dimensionnellement ; l'unité semble
nue. Cette nudité est une vertu pour la trigonométrie, mais c'est un oubli
pour une théorie qui entend expliquer la structure du continu et l'origine
de l'action. Le radion récupère l'oublié sans modifier la valeur publiée.

La généalogie commence dans APP, où une sortie conserve la valeur et aussi
le feuillet, le quotient, le reste, la retenue et la frontière. Le TRIT introduit régime et
orientation. Le TPK fait revenir la phase sans réinitialiser l'état. La connexion
nonadique construit la section coinductive de \(\pi\), et la dodécaphase publie
la charnière de \(50^\circ\). Le couple angulaire apporte un mode commun \(A\) et une
ouverture orientée \(C^*\). L'incidence de six phases actives produit
\(90\), ses deux feuillets sur l'espace ambiant \(729\) produisent \(20/81\), et la
torsion globale apporte la seconde section.

Alors apparaît \(\varepsilon_R\). Non comme une décimale implorée par la
clôture, mais comme la différence que la structure était tenue de
transporter :
\[
\epsone=\frac{20}{81}\Delta_4-(S_E-1).
\]
La retenue publie ses chiffres ; la profondeur certifie sa limite ; la carte
angulaire le reconnaît comme
\[
5.13830167585752820716851646226459474899\ldots
\times10^{-8}\,{}^\circ.
\]
L'équation du radian est clôturée parce que les deux membres sont des publications du
même état :
\BigRadionEquation

L'histoire ne s'arrête pas à l'égalité. Le couple \((A,C^*)\) construit une
ellipse positive. Son produit de demi-axes fixe \(2\hret\) ; son quotient fixe la
forme. Le théorème d'aire démontre qu'une unité de phase balaie \(\hret\) et
que le tour enferme \(\hfull\). C'est ici que la division
par \(2\pi\) acquiert un contenu physique : \(h\) est l'action de tour ; \(\hbar\) est l'action par radion.

Les foyers cessent d'être des ornements d'une figure. Leur séparation reconstruit
\(C^*\) dans la carte spectrale, et leur distance de Klein mesure une profondeur
hyperbolique finie au sein d'un domaine dont l'aire hyperbolique totale est
infinie. La même frontière contient une action finie. Le paradoxe apparent
est la thèse centrale de la géométrie : une mesure clôt l'action ; une autre ouvre la
profondeur.

Le cercle n'est donc pas faux. C'est la projection exacte qui conserve la phase
et oublie l'anisotropie, l'intérieur et la mémoire. L'ellipse conserve l'action et la forme.
Klein conserve la profondeur. L'hélice conserve le retour et la biographie. Aucune
image ne remplace les autres ; toutes naissent du même état.

La réalisation de Hilbert reconnaît l'ellipse comme contour de covariance
minimale. L'oscillateur LC la reconnaît comme échange de réserves électrique et
magnétique. Le lecteur \(90/120\) publie perméabilité, permittivité,
impédance et vitesse. L'incidence six--huit produit un écran causal ;
Hodge convertit l'orientation en dualité des champs ; Lorentz convertit champ et
charge en dynamique ; Cartan--Holst reçoit l'holonomie en géométrie
gravitationnelle. L'ordre importe : l'objet HMT sélectionne ; la physique
conventionnelle réalise.

Le mot \emph{radion} se trouve ainsi justifié. Il ne signifie pas qu'un foyer soit
une charge positive et l'autre une charge négative. Il signifie que l'unité de
phase possède une orientation signée avant la publication électrique :
\[
\mathcal A(\mathfrak r_+)=+\hret,
\qquad
\mathcal A(\mathfrak r_-)=-\hret.
\]
La charge apparaît lorsque la mécanique dimensionnelle applique le transducteur. Le signe
n'est pas ajouté à la fin ; il est conservé depuis le feuillet.

On explique aussi pourquoi le secteur \(90/120\) semblait cacher de nombreuses
réponses. Incidence, extracteur dodécaphasique, queue spectrale, noyau de
Barbero, tours et retenue partagent une origine. C'est précisément parce qu'ils partagent une origine
qu'ils peuvent être comparés. C'est précisément parce qu'ils ont des opérations distinctes qu'ils ne doivent pas
être fusionnés. Le typage ne réduit pas l'unité du système : il la rend intelligible.

% REMATE_LOCAL_RADION_CIERRE
\Needspace{28\baselineskip}
La conclusion la plus compacte de cet ouvrage peut s'écrire en quatre lignes :
\[
\boxed{
\begin{aligned}
\text{phase :}\quad&\theta=1,\\
\text{action :}\quad&\mathcal A=\sigma\hret,\\
\text{mémoire :}\quad&H_9(n+9)=H_9(n)+(0,1),\\
\text{compatibilité :}\quad&1=S_E-\Phi_T+\epsone.
\end{aligned}}
\]

Le radion est cette unité quadruple. Le radian est son ombre exacte dans la carte
angulaire. La différence n'est pas dans la valeur de \(\pi\) ; elle réside dans la quantité de
structure que nous sommes disposés à conserver lorsque nous disons qu'un tour s'est
fermé.

\vspace{2\baselineskip}
\begin{center}
\begin{minipage}{.82\textwidth}
\centering\large\itshape\color{RadNavy}
Le cercle publie la phase.\\
L'ellipse contient l'action.\\
Klein ouvre la profondeur.\\
L'hélice se souvient du chemin.\\[.7em]
Le radion est l'unité qui n'oblige pas à choisir entre elles.
\end{minipage}
\end{center}
